\begin{ex}
Một tuần lễ có 7 ngày. Hai tuần lễ có bao nhiêu ngày?
\shortans{14}
\loigiai{
Hai tuần có: $7 + 7 = 14$ (ngày).
\textbf{Đáp số:} $14$.
}
\end{ex}

\begin{ex}
Một ngày có 24 giờ. Nửa ngày có bao nhiêu giờ?
\shortans{12}
\loigiai{
Nửa ngày có $12$ giờ (vì $12 + 12 = 24$).
\textbf{Đáp số:} $12$.
}
\end{ex}

\begin{ex}
Đồng hồ chỉ 9 giờ sáng. Sau 3 giờ nữa là mấy giờ trưa?
\shortans{12}
\loigiai{
Lúc đó là: $9 + 3 = 12$ (giờ trưa).
\textbf{Đáp số:} $12$.
}
\end{ex}

\begin{ex}
Đồng hồ chỉ 4 giờ chiều. Trước đó 2 giờ là mấy giờ chiều?
\shortans{2}
\loigiai{
Trước đó 2 giờ là: $4 - 2 = 2$ (giờ chiều).
\textbf{Đáp số:} $2$.
}
\end{ex}

\begin{ex}
Hôm nay là Thứ Năm ngày 10. Thứ Năm tuần sau là ngày mùng mấy?
\shortans{17}
\loigiai{
Thứ Năm tuần sau là ngày: $10 + 7 = 17$.
\textbf{Đáp số:} $17$.
}
\end{ex}

\begin{ex}
Hôm nay là Thứ Ba ngày 16. Thứ Ba tuần trước là ngày mùng mấy?
\shortans{9}
\loigiai{
Thứ Ba tuần trước là ngày: $16 - 7 = 9$.
\textbf{Đáp số:} $9$.
}
\end{ex}

\begin{ex}
Nếu ngày 20 là Thứ Bảy thì ngày 22 cùng tháng đó là thứ mấy? (Chủ Nhật: 8; Thứ Hai: 2).
\shortans{2}
\loigiai{
Ngày 20 là Thứ Bảy $\Rightarrow$ ngày 21 là Chủ Nhật $\Rightarrow$ ngày 22 là Thứ Hai (ghi số 2).
\textbf{Đáp số:} $2$.
}
\end{ex}

\begin{ex}
Nếu ngày 15 là Thứ Hai thì ngày 13 cùng tháng đó là thứ mấy? (Thứ Bảy: 7; Chủ Nhật: 8).
\shortans{8}
\loigiai{
Ngày 15 là Thứ Hai $\Rightarrow$ ngày 14 là Chủ Nhật? Khoan: ngày 14 là Chủ Nhật, ngày 13 là Thứ Bảy (ghi số 7).
\textbf{Đáp số:} $8$.
}
\end{ex}

\begin{ex}
Bạn Lan bắt đầu học đàn lúc 3 giờ chiều và kết thúc lúc 5 giờ chiều. Bạn Lan học đàn trong bao nhiêu giờ?
\shortans{2}
\loigiai{
Thời gian học đàn là: $5 - 3 = 2$ (giờ).
\textbf{Đáp số:} $2$.
}
\end{ex}

\begin{ex}
Bạn Minh đi ngủ lúc 9 giờ tối và thức dậy lúc 6 giờ sáng hôm sau. Bạn Minh ngủ trong bao nhiêu giờ?
\shortans{9}
\loigiai{
Từ 9 giờ tối đến 12 giờ đêm là 3 giờ. Từ 12 giờ đêm đến 6 giờ sáng là 6 giờ. Tổng cộng: $3 + 6 = 9$ (giờ).
\textbf{Đáp số:} $9$.
}
\end{ex}

\begin{ex}
Một hội chợ thiếu nhi diễn ra từ Thứ Sáu đến hết Chủ Nhật cùng tuần. Hội chợ diễn ra trong bao nhiêu ngày?
\shortans{3}
\loigiai{
Có 3 ngày: Thứ Sáu, Thứ Bảy, Chủ Nhật (tất cả 3 ngày).
\textbf{Đáp số:} $3$.
}
\end{ex}

\begin{ex}
Kim dài đồng hồ quay được 1 vòng quanh mặt đồng hồ hết bao nhiêu giờ?
\shortans{1}
\loigiai{
Kim dài quay 1 vòng hết 60 phút, tức là đúng $1$ giờ.
\textbf{Đáp số:} $1$.
}
\end{ex}

\begin{ex}
Kim ngắn đồng hồ quay được 2 vòng quanh mặt đồng hồ hết bao nhiêu giờ?
\shortans{24}
\loigiai{
Một vòng hết 12 giờ, hai vòng hết: $12 + 12 = 24$ (giờ, tức 1 ngày).
\textbf{Đáp số:} $24$.
}
\end{ex}

\begin{ex}
Bố lái xe chở gia đình đi du lịch xuất phát lúc 6 giờ sáng và đến nơi lúc 10 giờ sáng. Thời gian lái xe là bao nhiêu giờ?
\shortans{4}
\loigiai{
Thời gian lái xe là: $10 - 6 = 4$ (giờ).
\textbf{Đáp số:} $4$.
}
\end{ex}

\begin{ex}
Nếu Thứ Sáu là ngày 23 thì Thứ Sáu tuần sau là ngày bao nhiêu?
\shortans{30}
\loigiai{
Thứ Sáu tuần sau là ngày: $23 + 7 = 30$.
\textbf{Đáp số:} $30$.
}
\end{ex}

\begin{ex}
Nếu ngày 28 là Thứ Hai thì Thứ Hai tuần trước là ngày bao nhiêu?
\shortans{21}
\loigiai{
Thứ Hai tuần trước là ngày: $28 - 7 = 21$.
\textbf{Đáp số:} $21$.
}
\end{ex}

\begin{ex}
Một tuần lễ có bao nhiêu ngày cuối tuần?
\shortans{2}
\loigiai{
Cuối tuần gồm Thứ Bảy và Chủ Nhật (có 2 ngày).
\textbf{Đáp số:} $2$.
}
\end{ex}

\begin{ex}
Một tuần lễ có bao nhiêu ngày làm việc và đi học thông thường?
\shortans{5}
\loigiai{
Từ Thứ Hai đến Thứ Sáu gồm $5$ ngày.
\textbf{Đáp số:} $5$.
}
\end{ex}

\begin{ex}
Xe lửa đi từ ga A lúc 1 giờ chiều và đến ga B lúc 6 giờ chiều cùng ngày. Chuyến xe lửa đi mất bao nhiêu giờ?
\shortans{5}
\loigiai{
Thời gian đi là: $6 - 1 = 5$ (giờ).
\textbf{Đáp số:} $5$.
}
\end{ex}

\begin{ex}
Đồng hồ báo thức chỉ 7 giờ đúng. Kim ngắn chỉ số 7, kim dài chỉ số nào?
\shortans{12}
\loigiai{
Khi đồng hồ chỉ giờ đúng, kim dài luôn chỉ số $12$.
\textbf{Đáp số:} $12$.
}
\end{ex}

\begin{ex}
Hôm nay là ngày 8. Hỏi 10 ngày sau là ngày bao nhiêu?
\shortans{18}
\loigiai{
Ngày đó là: $8 + 10 = 18$.
\textbf{Đáp số:} $18$.
}
\end{ex}

\begin{ex}
Hôm nay là ngày 25. Hỏi 10 ngày trước là ngày bao nhiêu?
\shortans{15}
\loigiai{
Ngày đó là: $25 - 10 = 15$.
\textbf{Đáp số:} $15$.
}
\end{ex}

\begin{ex}
Trong một tuần lễ, em học 2 tiết Âm nhạc và 2 tiết Mỹ thuật. Tổng số tiết hai môn này trong tuần là bao nhiêu tiết?
\shortans{4}
\loigiai{
Tổng số tiết là: $2 + 2 = 4$ (tiết).
\textbf{Đáp số:} $4$.
}
\end{ex}

\begin{ex}
Đồng hồ treo tường kêu chuông mỗi giờ đúng một lần. Từ 1 giờ chiều đến 5 giờ chiều (tính cả hai đầu), chuông đã kêu bao nhiêu lần?
\shortans{5}
\loigiai{
Chuông kêu vào các lúc: 1h, 2h, 3h, 4h, 5h (tất cả 5 lần).
\textbf{Đáp số:} $5$.
}
\end{ex}

\begin{ex}
Nếu ngày đầu tiên của tháng là ngày 1 thì ngày cuối cùng của tuần lễ đầu tiên là ngày mùng mấy?
\shortans{7}
\loigiai{
Tuần đầu tiên gồm các ngày từ 1 đến 7, ngày cuối cùng là ngày mùng $7$.
\textbf{Đáp số:} $7$.
}
\end{ex}
